\section{Representing mathematical domains}

This section gives a reusable design guide.  The point is not to claim that all
listed deep theorems have thereby been proved.  It is to identify the exact
finite-query interfaces through which their constructive content can be
formalized.

\subsection{Natural, integer, rational, and finite algebra}

Naturals are native.  Integers may be represented as differences of naturals,
with equality witnessed by cross-addition.  The checked derivation library
already supplies a ring of differences for polynomial reasoning.  Rationals
are pairs \((a,b)\) with signed numerator and positive denominator; equality is
cross-multiplication, and reduced form is optional.

A finite group, ring, module, or category can be encoded by operation tables and
axiom checkers.  For a fixed presentation, tables may be computed rather than
certified.  For query-dependent finite structures, a master code stores the
tables, and bounded folds check associativity, identities, distributivity, or
other axioms.

\subsection{Polynomial algebra and linear algebra}

A polynomial of bounded degree is a coefficient array.  Multivariate monomials
use a fixed mixed-radix exponent code after the query specifies variable and
degree bounds.  Polynomial identities reduce to coefficient equalities or to
ring-normalized term identities.

A matrix is a flattened coefficient array.  Determinants, adjugates, traces,
and products are finite sums and products.  The existing linear-algebra module
already proves many such identities in the object theory.  \(\FQNF\) is useful
when a theorem supplies a matrix certificate: a rank minor, Smith reduction
trace, eigenvalue cluster, or finite linear-combination witness.

\subsection{Roots, algebraic numbers, and factorization}

An algebraic root observation should normally be an isolating interval or disc,
possibly containing a cluster with multiplicity.  A certificate contains:

\begin{itemize}
\item rational region data;
\item a root-count or sign-variation witness;
\item separation from excluded boundaries;
\item an error or diameter bound;
\item optional matching data between refinements.
\end{itemize}

At collisions, return clusters or symmetric functions.  Ordering individual
roots is a gap-sensitive operation.  A complete factorization of bounded degree
can be a coefficient-array certificate checked by multiplication; producing
one uniformly requires a separate constructive factorization theorem.

\subsection{Analysis}

Limits, series, integrals, derivatives, and analytic continuations are stated
with quantitative remainder data.

\begin{itemize}
\item A convergent series query asks for a partial sum and a tail bound at
precision \(p\).
\item An integral query asks for a rational partition, sample or enclosure data,
and an error certificate.
\item A derivative query asks for a finite difference or polynomial model with
a remainder bound on a rational neighborhood.
\item Analytic continuation uses a finite chain of overlapping rational discs,
Taylor coefficients, and Cauchy estimates.
\end{itemize}

The verifier is elementary once the constants and moduli are supplied.  The
hard theorem is often exactly the producer of those constants.

\subsection{Probability}

For finite rational laws, all expectations and moments are finite sums.  An
effective limit theorem query can include:

\begin{itemize}
\item truncation and tail parameters;
\item a rational threshold separated from atoms when a discontinuous CDF test
is used;
\item characteristic-function or moment error bounds;
\item a finite coupling or martingale-difference table.
\end{itemize}

General measure spaces are not needed inside iPRA if the theorem is formulated
through finite rational approximants, quantization moduli, and tail bounds.
Exact threshold claims need anti-concentration or a separation margin.

\subsection{Combinatorics and finite model theory}

Finite colorings, graphs, hypergraphs, permutations, tableaux, and proof traces
are direct single-code objects.  A counterexample or witness is selected by
bounded search.  A universal finite theorem is reflected by a bounded product
checker.  Ramsey-type existence arguments may be formalized by explicit
finite bounds; without a bound, qualitative finiteness is insufficient.

\subsection{Number theory and Diophantine equations}

Divisibility, gcd, modular arithmetic, bounded factorization, and finite descent
already have strong support in \code{IPRADerivations}.  A Diophantine search
query consists of a height bound and a polynomial checker.  Descent certificates
record each smaller solution and the strict measure decrease; the entire chain
may be a trace code.

For integral or rational points on a curve, the generic verifier can enumerate
and check all points below a supplied height.  The general theorems of Siegel or
Faltings require the producer of that height bound; \(\FQNF\) cannot and should
not manufacture it from bare finiteness.

\subsection{Algebraic geometry}

A practical finitary interface uses finite affine charts and witnessed
relations.

\begin{itemize}
\item An ideal is represented by finitely many generators.  Membership comes
with coefficient polynomials witnessing a combination; a bounded search needs
a degree/height envelope.
\item A localization fraction \(a/s\) is coded by numerator and denominator
terms.  Equality in the presence of zero divisors is witnessed by an exponent
\(n\) such that a suitable power kills the cross difference.
\item A principal-open cover carries a finite power-cover certificate.
\item Gluing data are finite tuples of local sections with overlap equalities
and a uniqueness proof.
\item A finite affine atlas and a requested Cech degree determine a finite
cochain complex.
\end{itemize}

The scheme-like object is therefore a protocol of finite chart and overlap
queries.  General quasi-compactness or effective Nullstellensatz bounds remain
producer theorems.

\subsection{Topology and manifolds}

A compact metric presentation supplies finite nets at requested scales and a
modulus relating names to net points.  A topological query may request:

\begin{itemize}
\item a finite simplicial or cubical complex at mesh \(2^{-p}\);
\item a carrier relation for a map;
\item a chain, cycle, boundary, or homology certificate;
\item a finite cover and its nerve;
\item a chart list with overlap and orientation data.
\end{itemize}

Nonunique simplicial selections are harmless when any two are contiguous; the
induced chain maps are then homotopic by an explicit finite prism operator.
The difficult theorem-specific work is to connect the finite complex to the
presented space with an explicit mesh/modulus argument.

\subsection{Category theory}

Category theory in iPRA should be read for fixed presentations and finite
queries, not as a category of all mathematical objects.  A finite category is
a bounded object/morphism table with reflected axioms.  A finite diagram,
cone, natural transformation, adjunction on a finite presentation, or finite
Kan extension is a bounded certificate problem.  Yoneda calculations reduce to
finite component equations.

For a fixed PR family of objects rather than a finite category, the query must
bound every part actually inspected.  A ``universal'' construction is proved by
an explicit mapping property over the bounded presented morphisms.

\subsection{Geometry and general relativity}

Coordinate geometry over rationals or coded algebraic extensions is exact.
Incidence, determinant, conic, and metric identities reduce to polynomial or
rational equations with denominator guards.  Differential-geometric claims are
formulated for finite symbolic expressions or quantitative chart data.

For Schwarzschild, Reissner--Nordstr\"om, Kerr, or Kerr--Newman calculations,
a finite query fixes the coordinate chart, rational parameter codes, tensor
component, and derivative order.  A certificate may consist of normalized
rational expressions and denominator nonvanishing conditions.  This supports
exact local tensor identities.  It does not by itself establish global
analytic extension, causal structure, or existence of real-valued coordinate
transformations without additional quantitative presentations.

\subsection{A unifying adapter schema}

Across all domains, an adapter should expose:

\begin{enumerate}
\item the packed query and its intended finite meaning;
\item the certificate code bound;
\item the fixed Boolean checker and reflected formula;
\item the producer theorem or explicit boundary premise;
\item the terminal extractor;
\item stability, congruence, and gap theorems;
\item the final declaration of type \(\Prf\,\PRAAx\).
\end{enumerate}

This schema is deliberately repetitive.  Repetition at the interface prevents
hidden nonconstructive assumptions in the implementation.
